\chapter{The Address-Validity Leak}
\label{chap:canon}

Speculative-exception fuzzing, armed with the source-agnostic detection of \autoref{chap:priming},
uncovers a previously unknown leak, the address-validity leak, that arises from a basic property
of address translation: the CPU's canonicality and alignment checks on loads, stores, and
indirect branches. This chapter explains the mechanism. Its reach across isolation boundaries and
the covert-channel demonstrations are in \autoref{chap:eval}.

\section{The oracle}
\label{sec:canon-oracle}

On AArch64, address translation accepts a virtual address only when the bits above the configured
virtual-address size are all zero (the TTBR0 range) or all one (the TTBR1 range). Any other value
is out of range and faults~\cite{arm_vmsa}. On the Neoverse~N3 and the Cortex-A715, a speculative
access to a canonical target allocates microarchitectural state as usual (an L1D line, a data or
instruction TLB entry, a BTB way), whereas an \emph{out-of-range} target is fast-rejected before it
allocates anything and leaves no footprint. The presence or absence of that footprint is thus a
one-bit oracle on the canonicality of the accessed pointer (\autoref{fig:canon-oracle}).

\begin{figure}[t]
\centering
\fitwidth{\begin{tikzpicture}[
  font=\footnotesize,
  ptr/.style={draw=black!55, rounded corners=2pt, align=center, inner sep=3pt,
              minimum width=30mm, minimum height=8mm, font=\ttfamily\footnotesize},
  acc/.style={draw=black!60, thick, rounded corners=2pt, fill=black!4,
              align=center, minimum width=24mm, minimum height=9mm, inner sep=2pt},
  foot/.style={draw=probe!60!black, thick, rounded corners=2pt, fill=probe!16,
               align=center, minimum width=26mm, minimum height=9mm, inner sep=2pt},
  nofoot/.style={draw=black!45, thick, rounded corners=2pt, fill=black!3, dashed,
               align=center, minimum width=26mm, minimum height=9mm, inner sep=2pt},
  ar/.style={-{Stealth[length=2mm]}, thick, draw=black!70},
  lbl/.style={font=\scriptsize, text=black!65, align=center},
]
  % in-range path (top)
  \node[ptr, fill=probe!12] (pin) at (0,1.1) {in-range\\ pointer};
  \node[acc] (ain) at (3.3,1.1) {speculative\\ access};
  \node[foot] (fin) at (7.0,1.1) {allocates a line\\ (L1D / TLB / BTB)};
  \node[lbl, anchor=west, text=probe!45!black] at (9.1,1.1) {footprint \textbf{present} $\Rightarrow$ oracle $=1$};
  \draw[ar] (pin) -- (ain);
  \draw[ar] (ain) -- node[lbl, above]{translates} (fin);
  % out-of-range path (bottom)
  \node[ptr, fill=culprit!10] (pout) at (0,-1.1) {out-of-range\\ pointer};
  \node[acc] (aout) at (3.3,-1.1) {speculative\\ access};
  \node[nofoot] (fout) at (7.0,-1.1) {fast-rejected\\ \emph{nothing allocated}};
  \node[lbl, anchor=west, text=culprit] at (9.1,-1.1) {footprint \textbf{absent} $\Rightarrow$ oracle $=0$};
  \draw[ar] (pout) -- (aout);
  \draw[ar] (aout) -- node[lbl, above, text=culprit]{\ding{55}\ reject} (fout);
  % tie-in
  \node[lbl, align=center, text=black!70] at (4.5,-2.5)
     {make a target's legality depend on a secret bit $\Rightarrow$ one footprint measurement answers one yes/no};
\end{tikzpicture}
}
\caption{The canonicality oracle. A speculative access to an in-range pointer allocates
microarchitectural state (a line, a TLB entry, a BTB way). An out-of-range access is fast-rejected
before it allocates anything. Making a target's legality depend on a secret bit turns the presence
or absence of that footprint into a one-bit read of the secret, and the illegal access only ever
happens on a squashed path, so the procedure never faults.}
\label{fig:canon-oracle}
\end{figure}

\section{Alignment}
\label{sec:canon-alignment}

The same fast-reject covers alignment. A speculative access that is misaligned for its access type
is rejected before it allocates. An SP-relative access requires 16-byte alignment and a
PC-relative or indirect-branch target requires 4-byte alignment, so the oracle equally reports
whether a computed address carries the correct low alignment bits.

\section{Transmitters}
\label{sec:canon-transmitters}

Because the check is a property of the address, it fires on every speculative operation that
consumes one. We verified three transmitter classes, the load, the store, and the indirect branch.
\autoref{tab:canon-transmitters} lists the encodings we tested and the microarchitectural state
each allocates on an in-range target. Every load and store variant we tried leaked, whereas the
software prefetch \reg{PRFM} did not.

\begin{table}[t]\centering\small
\setlength{\tabcolsep}{6pt}
\caption{Transmitters for the canonicality oracle, with the microarchitectural state each
allocates on an in-range target. \cmark/\xmark: leaked / did not leak.}
\label{tab:canon-transmitters}
\begin{tabular}{@{}l l l c@{}}
\toprule
Class & Encodings & Footprint & Leaks \\
\midrule
Load            & \reg{LDR}, \reg{LDRB}, \reg{LDP}, \reg{LDAR} & L1D, dTLB & \cmark \\
Store           & \reg{STR}, \reg{STRB}, \reg{STP}            & L1D, dTLB & \cmark \\
Indirect branch & \reg{BLR}, \reg{BR}                         & BTB, iTLB      & \cmark \\
Prefetch        & \reg{PRFM}                                  & none               & \xmark \\
\bottomrule
\end{tabular}
\end{table}

\section{Speculation windows}
\label{sec:canon-windows}

The out-of-range access must happen only on a transient path, and we reach it with three distinct
speculation primitives: a mistrained bounds check (Spectre-v1), a store-bypassed bound
(Spectre-v4), and a return shadow (Spectre-RSB). The stack-canary demonstration of
\autoref{sec:eval-exploit} drives a single gadget each of these three ways.

\section{The branch-target channel evades priming}
\label{sec:canon-l1d}

When the transmitter is an indirect branch, the leak rides the BTB, and we effectively read the
canonicality of an indirect-branch target. The BTB allocation itself is the leakage. There is no
separate footprint to probe within the same run, and in particular the L1D cache cannot capture it
within the run that caused it. To learn whether the speculative branch allocated an entry, one must
execute an indirect branch with the \emph{same} program counter again and observe the BTB
prediction. \revizor{} generates loop-free test cases, so that second execution cannot happen inside
a single input's run. The leak therefore leaves no mark on its own run's L1D trace, survives only in
the carried-over BTB state, and surfaces in a later run, which is the class of leaks priming
discards (\autoref{prop:priming-blind}). \revizor{}'s priming~\cite{revizor} therefore discards it,
while the source-agnostic detection of \autoref{chap:priming} reports it. Priming missed this leak
throughout a 48-hour campaign, whereas the source-agnostic detection finds it in $84$\,s on average
(\autoref{tab:eval-violations}).

\section{Portability}
\label{sec:canon-portability}

The cache and TLB channels are not exclusive to Arm, and they reproduce on our Intel Alder Lake
baseline as well. The BTB channel is microarchitecture-specific. It reproduces on the Neoverse~N3
and the Cortex-A715, but not on the Cortex-X3 or on Alder Lake (\autoref{tab:eval-violations}).
